% ============================================================================
% TECHNICAL APPENDIX 9A: Annotation Pipeline Tooling — Python Reference
% API integration, pipeline implementation, and tool comparison
% ============================================================================
% Source: /chapters/ch9/Ch9A_final.md
% ============================================================================

\chapter*{Technical Appendix 9A: Annotation Pipeline Tooling}
\addcontentsline{toc}{chapter}{Technical Appendix 9A: Annotation Pipeline Tooling}

\begin{chapterquote}{Formal foundations for Chapter 9}
Chapter~9 introduced the annotation economy and the quality control mechanisms that keep it from producing noise. This appendix provides Python reference implementations for production annotation pipelines using Label Studio, inter-annotator agreement measurement, and multi-stage quality routing.
\end{chapterquote}

% ============================================================================
\section{Label Studio API Integration}
% ============================================================================

Label Studio provides an open-source annotation platform with a Python SDK supporting programmatic project management, task import, and export.

\begin{lstlisting}[style=python, caption={Label Studio connection and project creation}]
from label_studio_sdk import Client

ls = Client(url='http://localhost:8080', api_key='YOUR_API_KEY')
ls.check_connection()

# Text classification project
project = ls.start_project(
    title='Document Classification',
    label_config='''
    <View>
      <Text name="text" value="$text"/>
      <Choices name="label" toName="text" choice="single">
        <Choice value="relevant"/>
        <Choice value="irrelevant"/>
        <Choice value="uncertain"/>
      </Choices>
    </View>
    '''
)
\end{lstlisting}

\subsection{Task Import with Pre-Annotations}

\begin{lstlisting}[style=python, caption={Importing tasks with model pre-annotations}]
# Import with pre-annotations (model-assisted annotation)
pre_annotated_tasks = [
    {
        "text": "Document text here.",
        "predictions": [{
            "result": [{
                "from_name": "label",
                "to_name":   "text",
                "type":      "choices",
                "value":     {"choices": ["relevant"]},
                "score":     0.87
            }]
        }]
    }
]
project.import_tasks(pre_annotated_tasks)
\end{lstlisting}

\subsection{Export and Parsing}

\begin{lstlisting}[style=python, caption={Annotation export and label parsing}]
import pandas as pd

annotations = project.export_tasks(export_type='JSON_MIN')
results = []
for task in annotations:
    if task.get('annotations'):
        ann   = task['annotations'][0]
        label = ann['result'][0]['value']['choices'][0]
        results.append({
            'id':         task['id'],
            'text':       task['data']['text'],
            'label':      label,
            'annotator':  ann.get('completed_by', {}).get('id'),
            'created_at': ann.get('created_at'),
        })
df = pd.DataFrame(results)
\end{lstlisting}

% ============================================================================
\section{Complete HITLAnnotationPipeline}
% ============================================================================

The following class integrates quality control, inter-annotator agreement measurement, and expert routing into a single pipeline object.

\begin{lstlisting}[style=python, caption={Production annotation pipeline with quality control}]
import numpy as np
from label_studio_sdk import Client
from sklearn.metrics import cohen_kappa_score
from collections import defaultdict
from typing import Dict, List, Tuple
import logging

logger = logging.getLogger(__name__)

class HITLAnnotationPipeline:
    """
    Production annotation pipeline with quality control.

    Features:
    - Automatic overlap sampling for IAA measurement
    - Per-annotator kappa computation with outlier detection
    - Multi-stage routing: fast pass -> expert review
    """

    def __init__(self, ls_url: str, api_key: str, project_id: int,
                 min_kappa: float = 0.60,
                 overlap_fraction: float = 0.10,
                 expert_threshold: float = 0.45):
        self.client   = Client(url=ls_url, api_key=api_key)
        self.project  = self.client.get_project(project_id)
        self.min_kappa        = min_kappa
        self.overlap_fraction = overlap_fraction
        self.expert_threshold = expert_threshold

    def sample_overlap_items(self, items: List[dict]
                              ) -> Tuple[List[dict], List[dict]]:
        n          = len(items)
        n_overlap  = max(30, int(n * self.overlap_fraction))
        np.random.shuffle(items)
        return items[:n_overlap], items[n_overlap:]

    def compute_project_kappa(self) -> Tuple[float, List[str]]:
        tasks = self.project.export_tasks(export_type='JSON')
        overlap = [t for t in tasks
                   if len(t.get('annotations', [])) >= 2]
        if not overlap:
            return float('nan'), []

        per_task: Dict[int, Dict[str, str]] = {}
        for task in overlap:
            per_task[task['id']] = {}
            for ann in task['annotations']:
                aid   = str(ann.get('completed_by', {}).get('id', 'unknown'))
                if ann.get('result'):
                    label = ann['result'][0]['value']['choices'][0]
                    per_task[task['id']][aid] = label

        all_ann = list({a for labels in per_task.values()
                        for a in labels})
        pair_kappas = {}
        for i in range(len(all_ann)):
            for j in range(i + 1, len(all_ann)):
                a, b = all_ann[i], all_ann[j]
                la, lb = [], []
                for labels in per_task.values():
                    if a in labels and b in labels:
                        la.append(labels[a]); lb.append(labels[b])
                if len(la) >= 5:
                    pair_kappas[(a, b)] = cohen_kappa_score(la, lb)

        if not pair_kappas:
            return float('nan'), []

        all_a, all_b = [], []
        for (a, b) in pair_kappas:
            for labels in per_task.values():
                if a in labels and b in labels:
                    all_a.append(labels[a]); all_b.append(labels[b])
        aggregate = cohen_kappa_score(all_a, all_b) if all_a \
                    else float('nan')

        avg_kappa = defaultdict(list)
        for (a, b), k in pair_kappas.items():
            avg_kappa[a].append(k); avg_kappa[b].append(k)
        low_quality = [aid for aid, ks in avg_kappa.items()
                       if np.mean(ks) < self.expert_threshold]

        return aggregate, low_quality

    def route_for_expert_review(
            self, completed: List[dict]
    ) -> Tuple[List[dict], List[dict]]:
        _, low_quality = self.compute_project_kappa()
        accepted, expert_queue = [], []
        for task in completed:
            anns = task.get('annotations', [])
            if len(anns) >= 2:
                labels = {ann['result'][0]['value']['choices'][0]
                          for ann in anns if ann.get('result')}
                if len(labels) > 1:
                    expert_queue.append(task); continue
            ann_ids = [str(ann.get('completed_by', {}).get('id', ''))
                       for ann in anns]
            if any(aid in low_quality for aid in ann_ids):
                expert_queue.append(task)
            else:
                accepted.append(task)
        return accepted, expert_queue

    def cost_estimate(self, n_items: int,
                      cost_per_annotation_usd: float,
                      n_annotators_per_overlap: int = 3) -> dict:
        n_overlap   = max(30, int(n_items * self.overlap_fraction))
        n_regular   = n_items - n_overlap
        total_anns  = n_regular + n_overlap * n_annotators_per_overlap
        total_cost  = total_anns * cost_per_annotation_usd
        return {
            'n_items':              n_items,
            'n_overlap_items':      n_overlap,
            'total_annotations':    total_anns,
            'total_cost_usd':       round(total_cost, 2),
            'qc_overhead_fraction': n_overlap * n_annotators_per_overlap
                                    / total_anns,
        }
\end{lstlisting}

% ============================================================================
\section{Tool Comparison}
% ============================================================================

\begin{table}[h]
\centering
\caption{Annotation platform comparison}
\begin{tabular}{lllll}
\hline
\textbf{Tool} & \textbf{Data types} & \textbf{Managed} & \textbf{Self-host} & \textbf{Best use case} \\
\hline
Label Studio & All           & No  & Yes      & Research, privacy-sensitive \\
Argilla      & NLP, LLM eval & No  & Yes      & RLHF preference collection \\
CVAT         & Image, video  & No  & Yes      & Object detection, segmentation \\
Prodigy      & NLP           & No  & License  & Fast NLP with active learning \\
Scale AI     & All           & Yes & No       & High-volume commercial \\
SageMaker GT & All           & Yes & AWS      & AWS-integrated workflows \\
Labelbox     & All           & Yes & No       & Enterprise, collaborative \\
\hline
\end{tabular}
\end{table}

% ============================================================================
\section{Exercises}
% ============================================================================

\begin{Exercise}[Label Studio Setup]
Using the Label Studio SDK, set up a binary text classification project. Import 50 sample items. Annotate 20 items manually. Export the results and compute the distribution of labels.
\end{Exercise}

\begin{Exercise}[Majority Vote Resolution]
Extend the \texttt{HITLAnnotationPipeline} class to support configurable majority vote resolution for overlap items. When 3~annotators label an item and 2~agree, accept the majority. When all~3 disagree on a 3-class task, route to expert review.
\end{Exercise}

\begin{Exercise}[Real-World IAA Analysis]
Using the SNLI or MultiNLI datasets from HuggingFace (which include per-annotator labels), compute pairwise kappa for all annotator pairs. Identify the 10\% most contested items. What is the kappa improvement from removing them?
\end{Exercise}

\begin{Exercise}[Cost Estimation]
Using the \texttt{cost\_estimate} function, calculate the total cost for a project of 10{,}000 items at \$0.10/annotation with 10\% expert re-annotation at \$1.00/annotation. How does the optimal overlap fraction change if expert re-annotation is 20\% of items?
\end{Exercise}

\bigskip
\begin{center}
\rule{0.5\textwidth}{0.4pt}
\end{center}
\bigskip

\noindent\textit{This appendix provides production-ready Python tooling for annotation pipelines, inter-annotator agreement measurement, and quality-controlled expert routing, as introduced in Chapter~9.}
